%================================================================
% 作業ログ 2026-07-07（発表会前日・最終調整）
%   音源差し替え，Processing バグ修正，ゲームモード修正，
%   音質改善，MOP 検証．
%
%   ビルド:
%     cd work/shiozawa/work-0707/作業ログ0707
%     docker run --rm -v "$(pwd):/workspace" -w /workspace \
%       ghcr.io/paperist/texlive-ja:debian latexmk 作業ログ0707.tex
%================================================================
\documentclass[uplatex,dvipdfmx,a4paper,11pt]{jsarticle}

% --- 余白 ---
\usepackage[top=25mm,bottom=25mm,left=25mm,right=25mm]{geometry}

% --- 表・箇条書き・図・色 ---
\usepackage{array}
\usepackage{tabularx}
\usepackage{booktabs}
\usepackage{enumitem}
\usepackage{graphicx}
\usepackage{url}
\usepackage[table]{xcolor}
\usepackage{amssymb}

% --- 段落間隔（Wordのような見た目に寄せる）---
\setlength{\parindent}{0pt}
\setlength{\parskip}{0.6em}

% --- セクション見出しを「番号. 太字」+ 章末に水平線を入れる ---
\usepackage{titlesec}
\titleformat{\section}
  {\normalfont\Large\bfseries}{\thesection.}{0.6em}{}
\titleformat{\subsection}
  {\normalfont\large\bfseries}{\thesubsection}{0.6em}{}
\titlespacing*{\section}{0pt}{1.2em}{0.6em}
\titlespacing*{\subsection}{0pt}{1.0em}{0.4em}

% --- 章末水平線（手で \sectionrule と書いて区切る）---
\newcommand{\sectionrule}{\par\medskip\noindent\rule{\linewidth}{0.6pt}\par\medskip}

% --- 箇条書きの記号を Word 風（・→○→数字）に揃える ---
\renewcommand{\labelitemi}{\textbullet}
\renewcommand{\labelitemii}{\textopenbullet}
\setlist[itemize]{leftmargin=1.6em,itemsep=0.1em,topsep=0.2em}
\setlist[enumerate]{leftmargin=2.0em,itemsep=0.1em,topsep=0.2em}

% \textopenbullet が無い環境向けのフォールバック
\providecommand{\textopenbullet}{\ensuremath{\circ}}

% --- 進捗ガントチャート用（グリッド表）---
\definecolor{barDone}{gray}{0.55} % 完了
\definecolor{barProg}{gray}{0.78} % 進行中
\definecolor{barPlan}{gray}{0.90} % 予定
\definecolor{nowCol}{gray}{0.30}  % 今週の見出し
\newcolumntype{G}{>{\centering\arraybackslash}p{2.0em}}
\newcommand{\bD}{\cellcolor{barDone}}
\newcommand{\bP}{\cellcolor{barProg}}
\newcommand{\bQ}{\cellcolor{barPlan}}

\begin{document}

%================================================================
% ヘッダ
%================================================================
{\LARGE\bfseries 作業ログ}\par\bigskip

報告者：塩澤 匠生

報告日：2026年7月7日

プロジェクト名：タクトーン~ IMUジェスチャーで奏でるArduinoオーケストラ~ \\

\sectionrule

%================================================================
% 1. 進捗サマリー
%================================================================
\section{進捗サマリー（要約）}

\begin{itemize}
  \item 全体進捗率：95\%（先週 90\% → 95\%．発表会前日の最終調整として
        音源差し替え・バグ修正・音質改善・MOP 検証を実施）
  \item 進捗状況（今週の変化）：
    \begin{itemize}
      \item 本番音源をチューバ/トロンボーンからフルート/オルガンに差し替え
      \item Processing の画面描画バグ（Serial.list() 例外）を修正
      \item 横振りによるゲームモード選択が 2 回目以降動作しないバグを修正
      \item 音色の ADSR・倍音エンベロープ・エフェクトを改善し，フルートを 1 オクターブ上げ
      \item MOP 検証を実施（検証プログラムによる評価）
    \end{itemize}
  \item 今週の主要成果：
    \begin{itemize}
      \item 発表会デモに向けた音源・音質の最終調整が完了
      \item 実機で発生していたクラッシュ・動作不良を解消
      \item MOP 検証で品質指標の達成状況を確認
    \end{itemize}
  \item 重大課題（影響度：なし）：
    \begin{itemize}
      \item {}[明日 7/8 の発表会に向け，準備完了]
    \end{itemize}
\end{itemize}

\sectionrule

%================================================================
% 2. 進捗ガントチャート（計画と現在地）
%================================================================
\section{進捗ガントチャート（計画と現在地）}

チームのガントチャート（\texttt{meetings/ガントチャーver1ト.xlsx}）のフェーズ区分（110 基本設計／
120 詳細設計／200 実装／300 テスト／400 ストレッチ）に，現在の進み具合を重ねたもの．
「今」の縦位置は本報告時点（7/7＝発表会前日）．

\medskip
{\footnotesize
\setlength{\tabcolsep}{2pt}
\renewcommand{\arraystretch}{1.2}
\begin{tabular}{@{}p{8.5em}|*{12}{G}@{}}
 & & & & & & & & & & & & \multicolumn{1}{c}{$\downarrow$今} \\
作業項目（フェーズ）
 & \cellcolor{barPlan}\tiny 4/15 & \cellcolor{barPlan}\tiny 4/22 & \cellcolor{barPlan}\tiny 4/29
 & \cellcolor{barPlan}\tiny 5/6 & \cellcolor{barPlan}\tiny 5/13 & \cellcolor{barPlan}\tiny 5/20
 & \cellcolor{barPlan}\tiny 5/27 & \cellcolor{barPlan}\tiny 6/3 & \cellcolor{barPlan}\tiny 6/10
 & \cellcolor{barPlan}\tiny 6/17 & \cellcolor{barPlan}\tiny 6/24 & \cellcolor{nowCol}\tiny\textcolor{white}{7/8} \\
\midrule
企画・テーマ決定        & \bD & \bD &     &     &     &     &     &     &     &     &     &     \\
110 基本設計            &     &     & \bD & \bD &     &     &     &     &     &     &     &     \\
120 詳細設計            &     &     &     &     & \bD &     &     &     &     &     &     &     \\
\textbf{200 実装}       &     &     &     &     &     &     &     &     &     &     &     &     \\
\quad test\_v2 輪唱・同期 &     &     &     & \bD & \bD & \bD & \bD &     &     &     &     &     \\
\quad test\_v3 ゲーム(firmware) &  &     &     &     &     & \bD & \bD & \bD & \bD &     &     &     \\
\quad test\_v3 Processing &   &     &     &     &     &     & \bD & \bD & \bD &     &     &     \\
\quad production 構築    &     &     &     &     &     &     &     &     &     &     & \bD & \bD \\
\textbf{300 テスト}     &     &     &     &     &     &     &     &     &     &     &     &     \\
\quad 単体・結合・システム &   &     &     &     &     &     &     &     & \bD & \bD & \bD &     \\
\quad 評価（精度・遅延） &     &     &     &     &     &     &     &     &     &     & \bD & \bD \\
400 ストレッチ（強弱ほか） &  &     &     &     &     &     &     & \bQ & \bQ &     &     &     \\
発表会                  &     &     &     &     &     &     &     &     &     &     &     & \multicolumn{1}{c}{$\bigstar$} \\
\bottomrule
\end{tabular}
}

\medskip
凡例：\colorbox{barDone}{~~}~完了 \quad \colorbox{barProg}{~~}~進行中 \quad
\colorbox{barPlan}{~~}~予定 \quad $\bigstar$~発表会（7/8）

\medskip
※ 今週は発表会前日の最終調整として，音源差し替え・バグ修正・音質改善・MOP 検証を実施．
明日 7/8 の発表会に向けて準備完了の状態．

\sectionrule

%================================================================
% 3. 作業ログ（詳細トラッキング）
%================================================================
\section{作業ログ（詳細トラッキング）}

\begin{tabularx}{\linewidth}{@{}p{10em}p{6em}p{4em}p{3em}X@{}}
\toprule
作業項目 & 実施日 & 作業時間 & 進捗 & コメント \\
\midrule
音源差し替え & 2026-07-07 & 20 分 & 100\% & チューバ/トロンボーンの音源をフルート/オルガンに差し替え．音色 JSON の入替と AudioManager の読込パス修正 \\
Processing バグ修正 & 2026-07-07 & 40 分 & 100\% & 本番用 Processing の画面が描画されない問題を修正．Serial.list() の例外保護とポート検出のフォールバック処理を追加 \\
ゲームモード選択修正 & 2026-07-07 & 20 分 & 100\% & 横振りによるゲームモード選択が 2 回目以降動作しないバグを修正（applyPattern.cpp のナビゲート状態リセット） \\
音質改善 & 2026-07-07 & 40 分 & 100\% & 音色の ADSR・倍音エンベロープ・エフェクトを改善．フルートを 1 オクターブ上げて C5 に変更し，音域の分離を改善 \\
\midrule
\multicolumn{2}{@{}l}{\textbf{ソフトウェア 小計}} & \textbf{120 分} & & \\
\midrule
MOP 検証 & 2026-07-07 & 120 分 & 100\% & 検証プログラム（serial\_logger.py, analyze.py）を用いて MOP 項目の達成状況を評価．実機ログを収集し PASS/FAIL を判定 \\
\midrule
\multicolumn{2}{@{}l}{\textbf{合計}} & \textbf{240 分} & & \\
\bottomrule
\end{tabularx}

\sectionrule

%================================================================
% 4. 課題管理
%================================================================
\section{課題管理}

\begin{tabularx}{\linewidth}{@{}p{8em}Xp{2.5em}p{2.5em}Xp{3em}@{}}
\toprule
課題 & 発生原因 & 影響度 & 優先度 & 対策 & 状態 \\
\midrule
Processing 画面描画不可 & Serial.list() が環境依存で例外を投げる & 高 & 高 & 例外保護＋ポート検出フォールバックを追加 & 解決済 \\
ゲームモード 2 回目不可 & ナビゲート発火後の状態リセット漏れ & 中 & 高 & applyPattern.cpp で発火後に強制 Idle 復帰 & 解決済 \\
\bottomrule
\end{tabularx}

\medskip
※影響度＝プロジェクトへのインパクト，優先度＝対応の緊急性

\sectionrule

%================================================================
% 5. AI利用状況と検証
%================================================================
\section{AI利用状況と検証(再現性重視)}

\subsection{利用内容}

\begin{itemize}
  \item Claude Code を用いて Processing のバグ修正・音源差し替え・音質改善・
        ゲームモード選択のバグ修正を実施
  \item 具体的には SerialCore.pde の例外保護，AudioManager.pde の音色パラメータ調整，
        applyPattern.cpp のナビゲート状態リセット処理の生成を支援
  \item 生成コードは \texttt{pio run} でビルド成功を確認し，実機テストで動作検証済み
\end{itemize}

\sectionrule

%================================================================
% 6. 次回計画
%================================================================
\section{次回計画}

\begin{itemize}
  \item 目標：
    \begin{itemize}
      \item {}[7/8 発表会でデモを成功させる]
      \item {}[振り返りレポートの準備を開始する]
    \end{itemize}
  \item 達成基準：
    \begin{itemize}
      \item {}[発表会で指揮→楽器→PC の一連の流れが正常に動作する]
      \item {}[ゲームモードの選択・採点が正しく機能する]
    \end{itemize}
  \item 期限：
    \begin{itemize}
      \item {}[7/8 発表会当日]
    \end{itemize}
\end{itemize}

\sectionrule

%================================================================
% 7. その他
%================================================================
\section{その他}

\begin{itemize}
  \item {}[今週は発表会前日の最終調整に集中した．
        音源をチューバ/トロンボーンからフルート/オルガンに差し替えたのは，
        実機テストでの聴感評価を踏まえた判断]
  \item {}[Processing の Serial.list() が環境依存で例外を投げ，
        画面が描画されないクリティカルなバグを発見・修正した．
        ポート検出のフォールバック処理も追加し，異なる PC 環境でも
        安定して起動するようにした]
  \item {}[音色の ADSR・倍音エンベロープ・エフェクトを改善し，
        フルートを 1 オクターブ上げて C5 に変更することで，
        金管との音域分離を改善した]
  \item {}[MOP 検証を 2 時間かけて実施し，検証プログラムで
        各項目の達成状況を評価した]
  \item {}[明日 7/8 が発表会．本番環境は全バグ修正済みで，
        デモ可能な状態]
\end{itemize}

\sectionrule

%================================================================
% 8. 付録
%================================================================
\section{付録}

\begin{itemize}
  \item 添付資料：
    \begin{itemize}
      \item GitHub リポジトリ：\url{https://github.com/takushio2525/2026-hackathon-team23}
    \end{itemize}
  \item 主要コミット（塩澤担当分）：
    \begin{itemize}
      \item \texttt{f2c0757} --- チューバ・トロンボーンの音源をフルート・オルガンに差し替え
      \item \texttt{e7924ab} --- 本番用 Processing の画面が描画されない問題を修正
      \item \texttt{3e674c7} --- Processing のシリアルポート検出が動作しない問題を修正
      \item \texttt{63f9044} --- 横振りによるゲームモード選択が 2 回目以降動作しない問題を修正
      \item \texttt{f88a4ab} --- フルートを 1 オクターブ上げて C5 に変更
      \item \texttt{c84327a} --- 音色の ADSR・倍音エンベロープ・エフェクトを改善
    \end{itemize}
\end{itemize}

\sectionrule

\end{document}
